\section{Background}
\label{sec:background}

\subsection{LLM and Prompts}

Large Language Models (LLMs) are Transformer-based models trained to predict the next token $t_n$ in a sequence $T = (t_1, t_2, \dots, t_n)$ auto-regressively.
Trained on massive text corpora, they acquire broad linguistic and reasoning capabilities across NLP tasks such as generation, translation, summarization, and question answering.

An LLM is driven by a prompt $P$, typically composed of a system prompt $P_{sys}$ that specifies the model's role and constraints, and a user prompt $P_{user}$ that carries the request.
Since the output depends entirely on $P$, its content and structure directly determine the response, the very property that jailbreak attacks exploit.

\subsection{Safety Alignment}

To reduce misuse (e.g., generating misinformation or facilitating illegal activities), providers apply safety alignment so that models refuse harmful requests and act in line with human values and usage policies~\cite{instructgpt}.
Common techniques include supervised fine-tuning (SFT) and reinforcement learning from human feedback (RLHF), while more recent work enforces safety within the model's reasoning path through deliberative alignment~\cite{deliberative} and chain-of-thought~\cite{CoT}.

\subsection{Jailbreak Attack}

Despite these measures, LLMs remain vulnerable to jailbreak attacks, in which adversaries craft inputs that bypass alignment guardrails and elicit harmful responses~\cite{jailbroken}.
Attacks are categorized by the adversary's access: a white-box adversary has the parameters and gradients needed for optimization-based attacks, whereas a black-box adversary interacts only through the API.
Early research focused on single-turn attacks, where the entire malicious intent sits in one prompt, making it conspicuous to input filters and leaving no room to steer the model gradually.
Multi-turn attacks instead distribute the intent across conversational turns: each turn appears benign on its own, yet the accumulated context steers the model toward a harmful response while evading per-turn safety checks, a strategy that is structurally infeasible in the single-turn setting.

\subsection{Thompson Sampling}

Thompson sampling (TS) is a sequential decision-making algorithm that balances exploitation of current knowledge against exploration of new actions~\cite{thompson}.
It models each action's expected reward as a posterior distribution; at each step it draws one sample per action, selects the action with the highest sampled value, and updates that action's posterior from the observed reward.
Actions with greater uncertainty have wider posteriors and are thus explored more often, while consistently rewarding actions are increasingly favored as their posteriors concentrate.
In IAM-Teamer, we treat each attack strategy as an action and its judged success rate against a target as the unknown reward, using TS to rank strategies and surface the most promising as a hint each turn, while the attacker makes the final choice (Section~\ref{sec:strategy}).
